% ============================================================================
% ABSTRACT - REVISED
% ============================================================================

\begin{abstract}
Constraint Satisfaction Problems (CSPs) are fundamental to numerous real-world applications spanning task planning, resource allocation, and combinatorial optimization. Despite their impressive advances in language understanding and mathematical reasoning, Large Language Models struggle with structured constraint reasoning. Recent systematic evaluations reveal a "complexity cliff": as puzzle difficulty increases, LLM accuracy degrades sharply, approaching random performance when the underlying solver faces more than 20 conflicts. Current neural-symbolic approaches employ a translate-execute-repair pipeline where LLMs formalize constraints and repair translation errors when solvers fail. However, this paradigm exhibits two critical limitations: (1)~repair mechanisms lack formal semantic guidance, causing iterative cycling and stagnation; (2)~each problem is solved independently, discarding accumulated experience.

We introduce \textbf{PRISM} (Paradigm-guided Reasoning with Iterative Solver Memory), a memory-augmented neuro-symbolic framework that addresses both limitations through formally-grounded experience accumulation. PRISM operates in two stages: (1)~\textit{Offline}, the system distills verified solving paradigms from successful trajectories, where each paradigm encodes Z3-checkable pre/post-conditions enabling mechanical verification; (2)~\textit{Online}, retrieved paradigms guide LLM inference while a typed repair trajectory memory implements Jaccard-based stagnation detection and cosine-based loop detection to trigger adaptive strategy escalation. Fundamentally, paradigms receive formal certification before storage—distinguishing PRISM from prior memory-based methods relying on unverified natural language heuristics.

Evaluation on the ZebraLogic benchmark (600 training puzzles, 900 test puzzles across 3×5 to 6×6 scales) demonstrates consistent improvements over seven baselines, including pure chain-of-thought, established neuro-symbolic methods (Logic-LM, ExpeL), and neural-symbolic baselines. Notably, PRISM achieves a 14.4 percentage-point improvement over the prior neuro-symbolic approach (50.4\% vs.~36.0\%), with the largest gains (13.5pp) on the most difficult 6×6 puzzles—validating our hypothesis that experience accumulation becomes increasingly valuable for harder problems. The approach compresses 1500+ trajectories into ~35 verified paradigms (40:1 ratio) and reduces average repair rounds by 40--50\% compared to baselines.
\end{abstract}
